% Generated draft: requires mathematical and source review.
\documentclass[11pt]{article}
\usepackage[margin=25mm]{geometry}
\usepackage{fontspec}
\setmainfont{Latin Modern Roman}
\usepackage{amsmath,amssymb,mathtools,unicode-math}
\setmathfont{Latin Modern Math}
\usepackage{enumitem,needspace,xurl,hyperref}
\setlength{\parindent}{0pt}
\setlength{\parskip}{5pt}
\setlength{\emergencystretch}{3em}
\newcommand{\N}{\mathbb{N}}
\newcommand{\Z}{\mathbb{Z}}
\newcommand{\Q}{\mathbb{Q}}
\newcommand{\R}{\mathbb{R}}
\newcommand{\C}{\mathbb{C}}
\newcommand{\F}{\mathbb{F}}
\newcommand{\A}{\mathbb{A}}
\newcommand{\PP}{\mathbb{P}}
\newcommand{\E}{\mathbb{E}}
\newcommand{\eps}{\varepsilon}
\newcommand{\dd}{\,\mathrm{d}}
\begin{document}
\section*{UnsolvedMath ID: 10900043}
\textbf{Source problem code:} catalogue entry
\subsection{Definitions and mathematical statement}
Let $\Gamma\subset\operatorname{PSL}_2(\C)$ be a Kleinian group. Its limit set $\Lambda(\Gamma)\subset\widehat\C=\partial\mathbb H^3$ is the set of accumulation points of an orbit $\Gamma x$, for $x\in\mathbb H^3$. A Sierpi\'nski carpet means a space homeomorphic to the standard compact planar carpet obtained by iteratively deleting middle open squares.

Walsh's item 5.18 asks two characterization questions:
\begin{enumerate}
\item For which $\Gamma$ is there a subset $K\subset\Lambda(\Gamma)$ homeomorphic to a Sierpi\'nski carpet?
\item For which $\Gamma$ does $\Lambda(\Gamma)$ contain a continuum?
\end{enumerate}
A continuum conventionally means a nonempty compact connected space. To make the second question nontrivial, the likely intended meaning is a nondegenerate continuum, containing at least two points; otherwise every nonempty limit set contains a singleton continuum. This convention is not explicit and must be checked.

Containment is weaker than asking that the entire limit set be a carpet or a continuum. The question does not require the contained set to be $\Gamma$-invariant or to arise as another subgroup's limit set, and it does not restrict $\Gamma$ to convex-cocompact or finitely generated groups.

The registry's nested assertion that a nonelementary Kleinian limit set is necessarily a continuum cannot be used as stated: nonelementary Schottky groups have Cantor limit sets. Reported classifications for particular carpet groups therefore do not by themselves settle these unrestricted containment questions.
\subsection{Short English statement}
Characterize Kleinian groups whose limit sets contain a Sierpi\'nski carpet, and those whose limit sets contain a nontrivial compact connected subset, subject to clarifying the source's continuum convention.
\subsection{Sources}
\begin{itemize}
\item Local registry: \url{data/unsolvedmath/problems.json}, permanent UnsolvedMath ID $10900043$, exact source code \texttt{catalogue entry}, attributed to Walsh.
\item Delp, Hoffoss and Manning, \emph{Problems In Groups, Geometry, and Three-Manifolds} (2015), subsection 5.18, PDF page 5: \url{https://arxiv.org/pdf/1512.04620#page=5}; \url{https://arxiv.org/abs/1512.04620}. Supplied access date: 2026-07-29; source-tex extraction checked against rendered PDF text by the source record.
\item The registry mentions Kapovich--Kleiner classifications, Bonk--Kleiner constructions, and a work attributed to McMullen titled ``Kleinian groups with a carpet limit set,'' associated there with \url{https://arxiv.org/abs/math/0508227}. These are supplied attributions; bibliographic matching and exact hypotheses need verification.
\end{itemize}
\end{document}
